\documentclass{jsarticle}
\usepackage{amsmath} 
\usepackage{amssymb}
\begin{document}
\title{瞑想法と速読法、そして音読法}
\author{}
\date{}
\maketitle
  
 1:ボーとしている状態 
 2：何も考えていない状態
 3：頭頂葉を緩める状態
 4：側頭葉をリラックスする状態
 1－4の状態を維持する

 GSR速読法のジェネレティブな状態の頭頂葉の瞑想法が、
 頭頂葉を緩める仕方になっている。このリラックスで側頭葉も安定する。
 頭頂葉を緩められたら、ボーとして何も考えない。

 思考が浮かんできたら、それを沈めて消す。
 この行為が沈めることと消すことが快楽にもなっている。



 2005年のときにSRS速読法が出来た後に、2006年に胸と胃の部分に
 潜在意識が渦巻くので、御祓をしたり、縁側のドアを
 開けたり、散歩をしたりしていた。
 エビリファイからリスペリドンをやめたらこの症状が出て来た。
 リスペリドンを服薬したら3ヶ月で治った。
 潜在意識が渦巻くと瞑想法を頑張っていた。

 最近になって気づいたが、人の欠点が浮かんできたら、
 先生達はこの症状を使って、瞑想法をさせていたのに気づいた。

 結局は、SRS速読法は別にいけないわけではないらしい。

 思考が浮かんできたら、それを沈めることも合わせて消すのが快楽であるが、

 3月13日にすぐに身につくのをやめさせられて、4月17日に沈めて消すのも待たされて、

 石川先生と河相先生に頼りたいが、思考を沈めて消す行為に早く快楽を
 持たせたいとも思う。

 瞑想法には2種類の型がある。

 1：思考を止めること
 2：ブレインストーミング

 この2のブレインストーミングは、抵抗していてはリピートになるが、
 思考を沈めて消すのは、出来ると楽ともおもう。

 瞑想法は始めは命がけだが、できると楽で精神安定にもなる。

 先生達が瞑想法を進めるのも納得した。体験しなかったら瞑想法は
 いいとは全然分からなかった。



\end{document}
 
